% ============================================================
%  GUION DE DEFENSA --- 15 minutos
%  Proyecto: Efficient Fine-Tuning of Small Open LLMs for
%            Spanish Sentiment Analysis (LoRA/QLoRA)
%
%  Para DECIRLO en voz alta: frases cortas, no prosa de paper.
%  Arco: problema -> método -> resultados -> límites honestos
%        -> contribución -> cierre.
%  W1 (RoBERTuito) y W4 (H200/T4) pre-emptadas dentro del guion.
%  Cifras trazadas al paper YA corregido (Tablas I--V, §III--VI).
%  Marcas [apoyo visual] = sugerencia; el brief de diseño se
%  deriva DESPUÉS, no aquí.
% ============================================================

\documentclass[11pt,a4paper]{article}

\usepackage[utf8]{inputenc}
\usepackage[T1]{fontenc}
\usepackage{lmodern}
\usepackage{amssymb}
\usepackage{geometry}
\geometry{margin=2.4cm}
\usepackage{xcolor}
\usepackage{enumitem}
\usepackage{hyperref}
\hypersetup{colorlinks=true, linkcolor=blue!50!black, urlcolor=blue!50!black}

% ---- Comandos de apoyo ----
\newcommand{\bloque}[2]{\medskip\noindent{\large\textbf{\textcolor{blue!45!black}{#1}}}\hfill{\small\textbf{[#2]}}\par
  \nopagebreak\rule{\linewidth}{0.4pt}\par\nopagebreak\smallskip}
\newcommand{\clave}[1]{\noindent\textbf{Mensaje clave:} \emph{#1}\par\smallskip}
\newcommand{\visual}[1]{\noindent{\footnotesize\textcolor{teal!60!black}{$\square$~[apoyo visual]: #1}}\par\smallskip}
\newcommand{\fuente}[1]{\par{\footnotesize\textcolor{gray}{Cifras: #1}}\par}
\newcommand{\preempt}[1]{\par\smallskip\noindent{\footnotesize\textcolor{red!55!black}{\textbf{$\triangleright$ Pre-empción #1}}}\par}

% Lista de frases para decir (un guion = una frase = un respiro)
\newenvironment{decir}{\begin{itemize}[label=\textbf{--},leftmargin=1.3em,itemsep=2.5pt,topsep=3pt]}{\end{itemize}}

\title{\textbf{Guion de defensa --- 15 minutos}\\[2pt]
\large Fine-Tuning eficiente (LoRA/QLoRA) de LLMs pequeños para análisis de sentimiento en español}
\author{Material de preparación --- Fase 3 (para decir en voz alta)}
\date{}

\begin{document}
\maketitle

\noindent\textbf{Cómo ensayar.} Cada bloque lleva su \emph{tiempo acumulado} (esquina derecha) y un
\emph{mensaje clave} de una frase: si te pierdes, vuelve a esa frase. Cada gui\'on (\textbf{--}) es
una frase corta, pensada para decirse de un respiro. Las marcas {\footnotesize\textcolor{teal!60!black}{$\square$~[apoyo visual]}}
son sugerencias; las {\footnotesize\textcolor{red!55!black}{\textbf{$\triangleright$ Pre-empción}}} son
los puntos donde \emph{tú} sacas la debilidad antes que el tribunal. Ritmo objetivo: $\approx$130--140
palabras/min. Presupuesto: \textbf{13:30--14:00} hablado (contenido $\approx$13:45, \emph{incluida la demo de cierre en vivo}). El resto
del cuarto de hora es margen para respiraciones, una aclaración a media charla y el alargue natural
del directo.

% ============================================================
\bloque{1. Problema}{00:00 $\rightarrow$ 01:45 \,·\, 1 min 45 s}
\clave{Adaptar un modelo de lenguaje a una tarea era caro; la pregunta de este trabajo es cuándo y a qué coste deja de serlo.}
\visual{Título + tres tuits de ejemplo etiquetados +/$\,$neutral$\,$/$-$.}
\begin{decir}
  \item Buenos días. Os hablo de cómo adaptar modelos de lenguaje al análisis de sentimiento en español\ldots sin pagar una GPU cara.
  \item El problema es fácil de enunciar: hay millones de tuits en español con opiniones, y queremos clasificarlos en positivo, negativo o neutral.
  \item Pero lo difícil no es solo acertar. Es acertar \emph{barato}.
  \item Adaptar un modelo grande a la antigua ---reescribiendo todos sus pesos--- exige hardware caro y muchos datos etiquetados. Para casi todo el mundo, eso es inviable.
  \item Han aparecido técnicas que prometen hacerlo con una fracción del coste: LoRA y QLoRA. Pero para el español quedan preguntas sin responder.
  \item Y de esas preguntas va este trabajo: ¿cuándo merece la pena afinar?, ¿a qué coste?, ¿y sirve para distintas variedades del español?
\end{decir}

% ============================================================
\bloque{2. Método / Propuesta}{01:45 $\rightarrow$ 04:00 \,·\, 2 min 15 s}
\clave{Un estudio empírico reproducible en hardware gratuito: LoRA/QLoRA sobre dos Qwen3 pequeños, frente a los encoders de referencia, en tres cortes.}
\visual{Diagrama: modelos (Qwen3 1.7B/4B) · datasets (Cardiff, InterTASS) · los 3 cortes A/B/C.}
\begin{decir}
  \item Mi enfoque es un estudio empírico y reproducible. Afino dos modelos abiertos y pequeños: Qwen3 de 1.700 y de 4.000 millones de parámetros, que caben en una GPU gratuita de 16 GB.
  \item Como referencia, los encoders clásicos del español ---BETO y XLM-RoBERTa--- afinados a la antigua.
  \item Los datos: tuits de Cardiff para las dos primeras preguntas; e InterTASS ---España, Costa Rica y Perú--- para la tercera.
  \item Mido con macro-F1: la media de la F1 de las tres clases. Las trata por igual, así que no se puede ignorar la clase difícil, el \emph{neutral}.
  \item Y lo organizo en tres cortes: A, ¿cuándo afinar gana al prompting?; B, ¿gana LoRA a los encoders contando el coste?; C, ¿un adaptador de una variedad sirve para otra?
  \item Dos garantías de rigor: los hiperparámetros de LoRA están \emph{fijos} en todos los experimentos, y el test se toca una sola vez, al final.
\end{decir}

% ============================================================
\bloque{3. Resultados}{04:00 $\rightarrow$ 08:45 \,·\, 4 min 45 s}

\noindent\textbf{3a. Corte A --- punto de cruce} \hfill {\small 04:00 $\rightarrow$ 05:15}
\clave{El punto en que afinar supera al prompting no es universal: depende del tamaño del modelo.}
\visual{Figura A: macro-F1 vs n (escala log), con el cruce marcado para 1.7B y 4B.}
\begin{decir}
  \item Primer resultado, el corte A.
  \item Comparo afinar con LoRA frente a solo ``pedirle'' la respuesta al modelo con ejemplos en el prompt.
  \item Para el modelo de 1.7B, afinar gana enseguida: con menos de \textbf{siete} ejemplos etiquetados ya supera a su mejor prompting.
  \item Pero con el de 4B la historia cambia. Como de fábrica ya lo hace mejor, necesita más datos: no cruza con firmeza hasta unos \textbf{cincuenta} ejemplos.
  \item La conclusión no es un número mágico. Es un mecanismo: cuanto mejor es el modelo sin afinar, más datos necesita el fine-tuning para superarlo.
\end{decir}
\fuente{Tabla~II --- 1.7B mejor prompting 0.560, LoRA 0.590 en $n{=}7$; 4B mejor prompting 0.652, LoRA 0.666 en $n{=}50$.}

\medskip\noindent\textbf{3b. Corte B --- frontera de Pareto} \hfill {\small 05:15 $\rightarrow$ 08:00}
\clave{LoRA domina la frontera calidad-coste: más calidad que los encoders entrenando una fracción de los parámetros.}
\visual{Figura B: Pareto calidad vs parámetros (log). LoRA arriba a la izquierda; RoBERTuito marcado fuera de la frontera.}
\begin{decir}
  \item Segundo resultado, el corte B, y el más visual.
  \item A datos completos, LoRA gana a los encoders. El 4B llega a \textbf{0,707} de macro-F1; el 1.7B, a \textbf{0,698}. BETO se queda en 0,661; XLM-RoBERTa, en 0,646.
  \item Y lo hace entrenando una fracción mínima: frente a los encoders, entre \textbf{17 y 43 veces} menos parámetros con el 1.7B, y entre \textbf{9 y 24 veces} menos con el 4B.
  \item Comparado con afinar el modelo entero, son \textbf{268 veces} menos parámetros\ldots con prácticamente la misma calidad. Afinar el modelo entero no compensa.
\end{decir}
\preempt{W1 --- RoBERTuito (una frase y avanzar)}
\begin{decir}
  \item Un apunte honesto: en la tabla hay un modelo de fábrica, RoBERTuito, con 0,758; lo reporto, pero queda \emph{fuera} de la frontera a propósito, porque está afinado en otro corpus y lo evalúo sin reentrenar ---no es comparable en igualdad de condiciones.
  \item Contra mis baselines controlados ---BETO y XLM-R, mismo protocolo--- LoRA gana. Y sigo.
\end{decir}
\begin{decir}
  \item Y una pieza más: QLoRA. Cuantizando a 4 bits, retiene entre el \textbf{97 y el 99 por ciento} de la calidad usando aproximadamente la mitad de memoria.
  \item Eso mete el modelo de 4B en una tarjeta de consumo de 8 GB.
\end{decir}
\fuente{Tabla~IV --- LoRA 4B 0.707 / 1.7B 0.698; BETO 0.661, XLM-R 0.646; 17--43$\times$ (1.7B), 9--24$\times$ (4B), 268$\times$ vs full-FT; QLoRA 4B 0.701 a 6.79 GB (vs 16.7), 1.7B 0.681 a 4.93 GB (vs 8.72).}

\medskip\noindent\textbf{3c. Corte C --- transferencia dialectal} \hfill {\small 08:00 $\rightarrow$ 08:45}
\clave{Un adaptador transfiere casi siempre sin coste entre variedades; la excepción, replicada, es apuntar al español de España.}
\visual{Figura C: los dos heatmaps 3$\times$3 (1.7B y 4B), diagonal resaltada. (Detalle 3$\times$3 $\rightarrow$ explorador, Vista~3, en Q\&A.)}
\begin{decir}
  \item Tercer resultado, el corte C, en un titular: un adaptador entrenado en una variedad del español sirve para otra casi sin coste.
  \item Entreno en una y evalúo en otra ---España, Costa Rica, Perú---: hacia Costa Rica y Perú es gratis; la única excepción es apuntar al español de España, con una penalización de entre \textbf{0,031 y 0,039} de F1.
  \item Y lo que le da peso no es el número, sino que la asimetría se \emph{repite} en los dos tamaños de modelo. El detalle 3$\times$3 lo tengo en el explorador, si preguntáis.
\end{decir}
\fuente{Tabla~III --- penalización a ES 0.031--0.039 (2.6--4.6$\sigma$), replicada en 1.7B y 4B; transferencia a CR/PE dentro del ruido.}

% ============================================================
\bloque{4. Límites honestos}{08:45 $\rightarrow$ 11:15 \,·\, 2 min 30 s}
\clave{Los límites están declarados de antemano; varios resultados ``negativos'' no son fallos, son la prueba de que el experimento medía lo que decía.}
\visual{Lista breve: dev-split · sin test formal · H200 vs T4 · un solo dataset.}
\begin{decir}
  \item Paso a los límites, porque la honestidad es parte del trabajo.
  \item Sobre la penalización del corte C: en la memoria soy explícito. La describo en múltiplos del ruido entre semillas ---hasta 4,6 veces--- pero eso es una medida de \emph{magnitud}, no un test estadístico. No vendo una significancia que no he calculado.
  \item El corte C, además, lo evalúo en el conjunto de desarrollo, no en el de test. ¿Por qué? Porque el test oficial estaba inaccesible: daba error 404.
  \item El de desarrollo nunca se usó para entrenar ni para seleccionar, así que no hay trampa; solo aviso de que no es comparable con resultados ajenos sobre el test.
\end{decir}
\preempt{W4 --- H200 vs T4 (una frase y avanzar)}
\begin{decir}
  \item Una aclaración sobre el hardware: el título habla de una T4 gratuita de 16 GB, pero corrí en una H200.
  \item No es contradicción: lo que sostiene la accesibilidad es la \emph{memoria pico}, entre \textbf{4,9 y 8,7 GB} ---las cifras del objetivo T4--- y eso cabe en 16 GB corra donde corra.
  \item La H200 solo paraleliz\'o el barrido; lo único que cambiaría en una T4 es el tiempo. Y avanzo.
\end{decir}
\begin{decir}
  \item Hay más límites declarados: un solo dataset en los cortes A y B, corpus pequeños en el C, y dominio Twitter. Ninguno invalida los hallazgos centrales.
  \item Y un apunte, del lado bueno: probé la robustez a ruido real ---tildes, emojis, abreviaturas, mayúsculas--- y el único colapso claro es de los encoders \emph{cased} ante el mayúsculado (BETO pierde \textbf{0,18} de macro-F1; los LoRA, 0,01). Es un límite de los \emph{baselines}, no del método ---y os lo enseño en vivo al cerrar.
\end{decir}
\fuente{§V y §VI --- VRAM pico 4.93--8.72 GB (Tabla~IV); test InterTASS HTTP~404; 2.6--4.6$\sigma$ = heurístico, no test formal. Robustez: mayúsculas BETO $-$0.184 / XLM-R $-$0.110 (\texttt{results/robustness\_results.csv}, Fig.~F).}

% ============================================================
\bloque{5. Contribución}{11:15 $\rightarrow$ 12:15 \,·\, 1 min}
\clave{La aportación no es un récord de F1, sino un mapa reproducible de cuándo, cuánto y a qué coste conviene cada opción.}
\visual{Tres contribuciones, una línea cada una; abajo, la palabra ``reproducible''.}
\begin{decir}
  \item ¿Qué aporto? Tres cosas y un hilo común.
  \item Primero, el punto de cruce prompting-a-fine-tuning para el español: depende del tamaño del modelo.
  \item Segundo, un benchmark calidad-coste controlado donde LoRA domina a los encoders en hardware gratuito.
  \item Tercero, una primera evidencia de asimetría dialectal en la transferencia, replicada en dos modelos.
  \item Y el hilo común: reproducibilidad ---código, configuraciones y particiones liberadas---. No reclamo el mejor F1 absoluto, sino un mapa útil para quien decide con poco presupuesto.
\end{decir}

% ============================================================
\bloque{6. Cierre en vivo --- demo}{12:15 $\rightarrow$ 13:45 \,·\, 1 min 30 s}
\clave{No es una tabla más: el adaptador clasifica en directo, y el colapso cased se ve en vivo ---los encoders voltean con el mayúsculado, los LoRA aguantan.}
\visual{Demo en vivo servida desde la DGX (H200) por túnel SSH: los 5 modelos a la vez. Slide de respaldo con el antes/después estático por si falla (Plan~B).}

\noindent{\footnotesize\textcolor{teal!60!black}{$\square$~\textbf{Preparación (antes de la charla):} demo \emph{ya arrancada y con warmup hecho} en la DGX, túnel SSH abierto y probado, navegador en \texttt{localhost:7860}, frase lista para teclear. El primer clic ya va a tiempo pleno ($\approx$0,1\,s los 5 modelos); no hay espera de carga.}}\par\smallskip

\begin{decir}
  \item Os acabo de decir que los encoders \emph{cased} se hunden con las mayúsculas. En vez de otra tabla, os lo enseño en vivo.
  \item Tecleo una frase corta y positiva, pero sutil: ``pues no está nada mal''.
  \item Los cinco modelos la clasifican a la vez ---LoRA de 1.7B y 4B, prompting, BETO y XLM-R---. Todos: \emph{positivo}. De acuerdo.
  \item Ahora pulso el botón de MAYÚSCULAS ---el mismo ruido de la robustez--- y vuelvo a clasificar la misma frase.
  \item Y mirad: BETO y XLM-R voltean a \emph{negativo}. Solo por las mayúsculas. Los LoRA no se mueven: siguen en positivo.
  \item Ahí está el hallazgo, en directo: el encoder \emph{cased} confunde el grito con enfado; el adaptador generativo, no.
  \item Cierro como empecé: para el español, afinar barato es viable ---y aguanta el ruido real---. Gracias; quedo a vuestra disposición para las preguntas.
\end{decir}

\noindent{\footnotesize\textcolor{red!55!black}{\textbf{$\triangleright$ Plan B (obligatorio, si la demo falla):}}}
\begin{decir}
  \item Si el túnel o la DGX no responden: arranco el \emph{fallback} en CPU en el portátil (\texttt{python demo/app.py -{}-cpu}) ---misma demo, LoRA-1.7B, $\approx$1\,s por frase---.
  \item Si tampoco: salto a la \emph{slide de respaldo} con la captura del antes/después (BETO y XLM-R positivo$\rightarrow$negativo; LoRA positivo) y lo cuento con la misma frase.
  \item Regla: nunca me quedo en blanco ---la slide de respaldo está en el deck (oculta) y el mensaje es el mismo con o sin directo.
\end{decir}
\fuente{Demo: \texttt{demo/app.py} (GPU, 5 modelos) / \texttt{-{}-cpu} (fallback CPU). Colapso cased trazado a \texttt{results/robustness\_results.csv} (BETO $-$0,184; XLM-R $-$0,110).}

% ============================================================
\bigskip\noindent\rule{\linewidth}{0.8pt}
\section*{Puente al Q\&A (10 min) --- enlace con el banco \texttt{02\_banco}}

\noindent Las preguntas más probables ya las he \emph{tocado} en el guion; en el turno de preguntas
solo tengo que \textbf{rematar con la versión de 10 segundos} del banco. Atajo mental:

\begin{itemize}[leftmargin=1.4em,itemsep=2.5pt]
  \item \textbf{RoBERTuito (0,758)} $\rightarrow$ ya pre-emptado en 3b: ``fuera de Pareto, otro corpus, sin reentrenar''. (W1)
  \item \textbf{H200 vs T4} $\rightarrow$ ya pre-emptado en bloque 4: ``VRAM pico = cifras del target; solo cambia el tiempo''. (W4)
  \item \textbf{¿Los $\sigma$ son significancia?} $\rightarrow$ ``magnitud relativa al ruido, no test; y replica en dos modelos''. (W3)
  \item \textbf{Una sola semilla en B} $\rightarrow$ ``$\pm$0,006 a datos completos; irrelevante frente a 4--6 puntos''. (W5)
  \item \textbf{¿El factor de parámetros?} $\rightarrow$ ``10.000$\times$ es de Hu et al.; lo mío es 17--43$\times$ / 9--24$\times$ / 268$\times$''. (W14)
  \item \textbf{¿Uso de IA?} $\rightarrow$ ``yo orquesté y validé con gates; el diario es la traza''. (CE37/RA5)
  \item \textbf{¿Y la robustez?} $\rightarrow$ ``el único colapso es \emph{cased} ante mayúsculas; LoRA aguanta; lo acabáis de ver en la demo''. (Sprint 6)
\end{itemize}

\noindent\textbf{Puertas que NO abrir} (no las saques tú; si te empujan, recoge y recoloca):
\begin{itemize}[leftmargin=1.4em,itemsep=2pt]
  \item No compares tiempos de entrenamiento T4 vs H200 (no los medí en T4).
  \item No digas que Qwen3 es ``mejor'' que Llama/Gemma (no lo medí): defiende los criterios.
  \item No defiendas el 10.000$\times$ como cifra propia: es de Hu et al.
  \item No afirmes causalidad dialectal (por qué ES es sensible): el dato es la asimetría, no su causa.
  \item No presentes la robustez como exhaustiva: son 7 perturbaciones sobre un test; el mayúsculado es el hallazgo, no un \emph{survey}.
\end{itemize}

\noindent\textbf{Regla de oro:} si una pregunta toca un límite, recondúcela a su \emph{encuadre}
(``está declarado a propósito porque\ldots'') en vez de improvisar. Respuestas largas en el banco;
en la sala, la de 10 segundos.

\medskip
\noindent\textbf{Herramientas de respaldo} (Q\&A, \emph{no} en el discurso):
\begin{itemize}[leftmargin=1.4em,itemsep=2.5pt]
  \item \textbf{Pareto / cruce / transferencia} $\rightarrow$ abrir \texttt{defense/deck\_assets/explorer.html} (offline, doble clic): Vista~1 (Pareto, Corte~B), Vista~2 (cruce F1-vs-n, Corte~A), Vista~3 (transferencia 3$\times$3, Corte~C).
  \item \textbf{Robustez / otras perturbaciones} (tildes, emojis, abreviaturas\ldots) $\rightarrow$ \texttt{figF} (heatmap de robustez, \texttt{results/figures/figF\_robustness.png}); tabla completa en \texttt{results/robustness\_results.csv}.
  \item \textbf{¿Otra frase en la demo?} $\rightarrow$ la demo sigue abierta; pruebo la que propongan (o el modo \texttt{-{}-cpu} en el portátil).
  \item \textbf{Trade-off de coste 4D} $\rightarrow$ \texttt{figE} (dispersión 4D, asset de reserva).
\end{itemize}

\end{document}
